\section{Simetria}

\begin{def:simetria}
\label{def:simetria}
Uma variável $X$ possui distribuição simétrica entorno de $x = m$ se, e somente se,
\[
    \forall t \in \real,\, P(X - m \le t) = 1 - P(X - m < -t)
\]
\end{def:simetria}

\begin{obs:simetria 1}
Se fizémos $t = -u$, então
\begin{equation}
\label{eq:equacao simetria}
    P(X - m \le -u)  = 1 - P(X - m < u)
\end{equation}
para todo $u \in \real$.
Logo, a expressão \ref{eq:equacao simetria} é uma formulação equivalente para
a definição \ref{def:simetria}.
\end{obs:simetria 1}

\begin{obs:simetria 2}
Se $X$ for variável contínua, então a definição de simetria equivale a
\begin{gather*}
	P(X - m \le t) = 1 - P(X - m \le -t) \\
	F_X(m+t) = 1 - F_X(m-t)
\end{gather*}
\end{obs:simetria 2}

\begin{prop:equivalencia simetria}
\label{prop:equivalencia simetria}
Se $X$ é uma variável com distribuição simétrica entorno de $m$, então
\[
    P(X \le m) = P(X \ge m ).
\]
\end{prop:equivalencia simetria}
\begin{proof}
Com efeito,
\begin{gather*}
P(X - m \le 0) = 1 - P(X - m < 0) \\
P(X \le m) = 1 - P(X < m) \\
P(X \le m) - [1 - P(X < m)] = 0 \\
P(X \le m) - P(X \ge m) = 0 \\
P(X \le m) = P(X \ge m) \\
\end{gather*}
\end{proof}

\begin{def:mediana}
A mediana $m$ da distribuição de uma variável aleatória $X$ satisfaz
\[
    P(X \ge m) \ge \frac 1 2 \quad \text{e} \quad P(X \le m) \ge \frac 1 2
\]
\end{def:mediana}

\begin{prop:simetria e mediana}
Se $X$ possui distribuição simétrica entorno de $m$, então $m$ é mediana da distribuição
de $X$.
\end{prop:simetria e mediana}
\begin{proof}
Sabemos que
\begin{gather*}
	P(X \le m) + P(X > m) = 1 \\
	P(X \le m) + P(X > m) + P(X = m) = 1 + P(X = m) \\
	P(X \le m) + P(X \ge m) = 1 + P(X = m) \\
\end{gather*}
mas pela Proposição \ref{prop:equivalencia simetria}
\[
    2P(X \le m) = 1 + P(X = m) \\
\]
Uma vez que $P(X = m) \ge 0$, então $1 + P(X = m) \ge 1$, e
\begin{gather*}
    2P(X \le m) \ge 1 \\
    P(X \le m) \ge \frac 1 2. \\
\end{gather*}
Mais uma vez, pela Proposição \ref{prop:equivalencia simetria},
\[
    P(X \ge m) \ge \frac 1 2 \\
\]
\end{proof}

\begin{thm:simetria funcao de probabilidade}
\label{thm:simetria funcao de probabilidade}
A variável $X$ ter simetria entorno de $m$ equivale a
\begin{enumerate}
\item Se $X$ for discreta, então
\[
    p_X(m+t) = p_X(m-t).
\]

\item Se $X$ for contínua, então
\[
    f_X(m+t) = f_X(m-t).
\]
\end{enumerate}
\end{thm:simetria funcao de probabilidade}
\begin{proof}
Suponha que $X$ seja simétrica.
Para provar 1, considere que
\[
    P(X - m \le t) = P(X - m = t) + P(X - m < t)
\]
mas, por simetria,
\[
    P(X - m \le t) = 1 - P(X - m < -t)
\]
de forma que
\[
    P(X - m = t) = 1 - [P(X - m < -t) + P(X - m < t)]
\]
Analogamente, por $\sigma$-aditividade e simetria, iremos concluir que,
\[
    P(X - m = -t) = 1 - [P(X - m < -t) + P(X - m < t)]
\]
Logo,
\[
    P(X - m = t) = P(X - m = -t)
\]
e, como $X$ é discreto, então
\[
    p_X(m+t) = p_X(m-t)
\]
Para provar a afirmação 2, observe que
\[
    F_X(m+t) = 1 - F_X(m-t),
\]
e como $f_X$ é contínua em $\Omega_X$, então fazendo $F_Y(y) = F_X(m+t)$ e $F_U(u) = F_X(m-t)$,
obtemos
\[
     \frac{d}{dt} F_Y(y) = - \frac{d}{dt} F_U(u).
\]
Já que $y = m+t$ e $u = m-t$, então pelo Teorema \ref{thm:fdp transformacao continua}
\[
     f_X(m+t) = - [-f_X(m-t)] = f_X(m-t).
\]

Agora suponha que as equações das afirmações 1 e 2 valem.
Partindo de 1, temos que
\[
    P(X-m \le -t) = \sum_{y \le -t} P(X-m=y)
\]
mas como para todo $t \in \real$ vale que $P(X-m=t) = P(X-m=-t)$, então
\[
    P(X-m \le -t) = \sum_{y \le -t} P(X-m=y) = \sum_{-y \ge t} P(X-m=-y) = P(X-m \ge t).
\]
Uma vez que $P(X-m \ge t) = 1 - P(X-m < t)$, então
\[
    P(X - m \ge t) = 1 - P(X - m < t).
\]
O que prova a simetria para o caso discreto.
Direcionando agora para o caso contínuo, considere $Y = X - m$, de modo que
$f_y(y) = f_X(y+m)$.
Com efeito
\[
    F_Y(Y \le t) = \int_{-\infty}^{t} f_X(m+y)\,dy = 1 - \int_{t}^{\infty} f(m+y)\,dy.
\]
mas supondo a afirmação 2, obtemos
\[
    F_Y(Y \le t) = 1 - \int_{t}^{\infty} f(m-y)\,dy.
\]
Fazendo a substituição $y = -u$ e invertendo os limites de integração ajustados à troca
de variável, veremos que
\[
    F_Y(Y \le t) = 1 - \int_{-\infty}^{-t} f(m+u)\,du
\]
Fazendo outra troca de variável do tipo $u = w$, é possível concluir que
\begin{gather*}
    F_Y(Y \le t) = 1 - F_Y(Y \le -t) \\
    F_X(m+t) = 1 - F_X(m-t)
\end{gather*}
\end{proof}

\begin{cor:simetria e media aritmetica}
Se $X$ tem distribuição simétrica entorno de $m$ e com função de probabilidade -- seja
de massa ou densidade -- $\pi(x)$, então
\[
    \forall x_1, x_2 \in \real,\, x_1 \le m \le x_2 \text{ e } \pi(x_1) = \pi(x_2) \rightarrow \frac{x_1 + x_2}{2} = m
\]
\end{cor:simetria e media aritmetica}
\begin{proof}
Suponha que $x_1 = m - a$ e $x_2 = m + b$, logo
\[
    \pi(x_1) = \pi(m-a) = \pi(m+a)
\]
mas, pela hipótese $\pi(x_1) = \pi(x_2)$, temos
\[
    \pi(m+a) = \pi(m+b).
\]
Logo, $a = b$ e, por conseguinte,
\[
    \frac{x_1 + x_2}{2}  = m.
\]
\end{proof}

\begin{def:moda}
A moda $\mo X$ de uma variável aleatória $X$ é o valor tal que
\begin{enumerate}
\item Se $X$ é discreta, então
\[
    \mo X = \argmax_x p_X(x)
\]
\item Se $X$ for contínua, então
\[
    \mo X = \argmax_x f_X(x)
\]
\end{enumerate}
\end{def:moda}

\begin{thm:simetria media moda mediana}
Se $X$ é uma variável aleatória cuja distribuição é simétrica entorno da mediana $m$,
então
\begin{enumerate}
\item Se a distribuição é unimodal, então $m = \mo X$.
\item Se $\ev{X} < \infty$, então $m = \ev X$.
\end{enumerate}
\end{thm:simetria media moda mediana}
\begin{proof}
Seja $p$ a função de massa de $X$ discreta.
Suponha que a moda esteja num valor $x$ diferente de $m$, de modo que $x = m + a$ ou
$x = m - a$.
No entanto, por hipótese de simetria, do Teorema \ref{thm:simetria funcao de probabilidade}
segue que
\[
    p(m+a) = p(m-a),
\]
implicando que a distribuição tem duas modas, logo não pode ser unimodal.
O mesmo argumento é replicado para uma variável contínua substituindo a função $p$ pela
densidade $f$.
Portanto, $m$ é a moda da distribuição.

Consideremos agoa a transformação $Y = X - m$, tal que
\begin{gather*}
P_X(Y \le t) = 1 - P_X(Y \le -t) \\
P_X(Y \le t) = P_X(Y \ge -t) \\
P_X(Y \le t) = P_X(-Y \le t),
\end{gather*}
indicando que $Y$ e $-Y$ tem a mesma distribuição, logo, o mesmo valor esperado.
Com efeito,
\begin{align*}
    \ev Y &= \ev{-Y} \\
    \ev{X} - m &= m - \ev X \\
    \ev X &= m.
\end{align*}
\end{proof}
